\documentclass[11pt,a4paper]{article}
\usepackage[margin=1in]{geometry}
\usepackage[hidelinks]{hyperref}
\usepackage{enumitem}
\usepackage{graphicx}
\usepackage{xcolor}

% -----------------------------------------------------
% bvraghav-tiet-ucs503p-article.sty
% -----------------------------------------------------

% Variable: Lab Instructor
\makeatletter \newcommand{\BvrSetLabInstructor}[1]{%
  \gdef\Bvr@LabInstructor{#1}}
% default
\newcommand{\Bvr@LabInstructor}{Ms. Paramveer %
  Kaur}
% public accessor
\newcommand{\BvrLabInstructorValue}{\Bvr@LabInstructor}
\makeatother

% Add subtitle
\usepackage{titling}
% -- Set title bold
\pretitle{\begin{center}\bfseries\fontsize{18bp}{18bp}\selectfont}
\makeatletter
\newcommand{\BvrSubtitle}[1]{%
  \posttitle{%
    \par\end{center}
    \begin{center}\large#1\end{center}
    \vskip 0.5em}%
}
\makeatother

% -----------------------------------------------------

\title{Scientific Literature Intelligence Platform}

\BvrSetLabInstructor{Ms. Paramveer Kaur}

\BvrSubtitle{%
  Project Proposal for UCS503P  \\
  Submitted to: \BvrLabInstructorValue \\ \bfseries
  Thapar Institute of Engineering and Technology%
}

\author{
  Paarth Ganesh, CSED%
  \thanks{Roll No: \texttt{1024030171}, %
    Email: \texttt{pganesh\_be24@thapar.edu}}
  \and
  Diya Jain, CSED%
  \thanks{Roll No: \texttt{1024031154}, %
    Email: \texttt{djain2\_be24@thapar.edu}}
}

\date{\today}

\begin{document}
\maketitle

\section{Higher-order goal}
This project builds a retrieval-augmented research assistant
that lets a user ask natural-language questions over a corpus
of scientific papers and receive answers that are grounded in,
and cited back to, specific source passages. The corpus is
populated both automatically---by fetching papers from the
arXiv API---and manually, by letting users upload their own
PDFs. The engineering objective is a deployed web service that,
given a question, returns an answer where every factual claim
is traceable to a retrieved chunk, plus per-paper and
cross-paper summaries. We treat retrieval quality and answer
faithfulness as first-class, measurable properties rather than
afterthoughts.

\section{Time-to-value}
\begin{itemize}[leftmargin=0.7cm]
    \item \textbf{Iteration 1 (end of week 2):} a working vertical slice---ingest a fixed set of arXiv PDFs, embed them, and serve semantic search over a REST endpoint. This alone is demonstrable and measurable.
    \item \textbf{Iteration 2 (week 4):} grounded question answering with inline citations, backed by a labeled evaluation set so every later change is scored, not guessed.
    \item \textbf{Iteration 3+ (week 5 onward):} summarization, structured knowledge extraction, hybrid retrieval, and reranking---each added only after it beats the previous version on the eval harness.
    \item CI runs the test suite and a retrieval-quality check on every commit, so regressions surface immediately rather than at demo time.
\end{itemize}

\section{Problem Statement}
Researchers, students, and analysts lose significant time reviewing literature across thousands of papers. The concrete failures we target are:
\begin{itemize}[leftmargin=0.7cm]
    \item \textbf{Information overload:} the volume of publications makes exhaustive manual review infeasible, so relevant prior work is missed.
    \item \textbf{Keyword-search limitations:} lexical search fails when authors describe the same concept with different terminology, hiding semantically related work.
    \item \textbf{Ungrounded answers:} general-purpose chatbots fabricate citations and DOIs, so their output cannot be trusted for factual, source-backed claims.
    \item \textbf{Slow synthesis:} manually extracting and comparing methods, datasets, and reported metrics across many papers is repetitive and error-prone.
\end{itemize}

\textbf{Why this matters now:} scientific output is growing faster than any individual can absorb, while retrieval-augmented generation has made grounded, citation-backed answering practical on commodity hardware. The platform converts that capability into saved review time and verifiable answers.

\section{Proposed Solution}
\subsection{Overview}
A \textbf{web-based} platform composed of five subsystems:
\begin{itemize}[leftmargin=0.7cm]
    \item \textbf{Ingestion:} fetch papers from the arXiv API (or accept uploaded PDFs), extract text and metadata, chunk, embed, and store.
    \item \textbf{Semantic search:} embedding-based retrieval that surfaces conceptually relevant passages rather than exact keyword matches.
    \item \textbf{Grounded question answering:} answers generated strictly from retrieved passages, each claim carrying an inline citation to its source paper and page.
    \item \textbf{Summarization:} concise summaries of a single paper or across a set of retrieved results.
    \item \textbf{Knowledge extraction:} structured pull of task, method, dataset, and headline metric, enabling side-by-side comparison across papers.
\end{itemize}

\subsection{Concrete technology stack}
Chosen for a cost-aware, industry-realistic build---local models on the high-volume retrieval path, a hosted LLM only for final generation.
\begin{itemize}[leftmargin=0.7cm]
    \item \textbf{Backend:} Python + FastAPI (async, typed, auto-generated OpenAPI docs).
    \item \textbf{Frontend:} React with a lightweight component library; chat-style Q\&A, search, and corpus views.
    \item \textbf{Embeddings (local):} \texttt{sentence-transformers} \texttt{all-MiniLM-L6-v2} (384-dim), run locally so retrieval has zero per-query API cost.
    \item \textbf{Vector store + database:} PostgreSQL with the \texttt{pgvector} extension---one system for both relational metadata and vector search, avoiding a second datastore.
    \item \textbf{Generation LLM (hosted):} Gemini 2.5 Flash, called only to compose the final grounded answer and summaries from retrieved context.
    \item \textbf{PDF parsing:} PyMuPDF for text and page-level positions (needed for page-accurate citations).
    \item \textbf{Ingestion source:} the official arXiv API for metadata and PDF links.
    \item \textbf{Deployment:} Docker Compose; app on Render/Railway, managed Postgres alongside.
    \item \textbf{CI/CD:} GitHub Actions---lint, tests, and a retrieval-quality gate on every push.
\end{itemize}

\subsection{Core workflow}
\begin{enumerate}[leftmargin=0.7cm]
    \item User fetches papers by arXiv query or uploads PDFs; the ingestion pipeline parses, chunks, embeds, and stores them.
    \item User poses a natural-language question.
    \item The retriever embeds the query and pulls the top-$k$ chunks by vector similarity, optionally reranked.
    \item The generator composes an answer strictly from those chunks, attaching a citation (paper + page) to each claim, and returns it with summaries and extracted fields.
\end{enumerate}

\section{Solution Approach}
This is an engineering project: the novelty is in system design, retrieval quality, and grounding discipline, not in inventing models.

\subsection{Ingestion and chunking}
\begin{itemize}[leftmargin=0.7cm]
    \item Parse PDFs with PyMuPDF, retaining page numbers so every chunk knows its origin.
    \item Chunk by a fixed token window ($\approx$500 tokens) with $\approx$15\% overlap, so context is not severed mid-idea.
    \item Store each chunk with its paper id, page, character span, and 384-dimensional embedding.
    \item Ingestion is idempotent: re-fetching a known arXiv id updates rather than duplicates it.
\end{itemize}

\subsection{Retrieval and grounding}
\begin{itemize}[leftmargin=0.7cm]
    \item Baseline: dense retrieval via cosine similarity on \texttt{pgvector} using an HNSW index for sub-linear search.
    \item Upgrade: hybrid retrieval---combine dense scores with lexical BM25, then rerank the merged top results with a cross-encoder for precision.
    \item Grounding contract: the generator receives only retrieved chunks and is instructed to answer solely from them; if the context does not support an answer, it says so rather than inventing one.
    \item Every claim in the output is tied to a chunk id, which resolves to a paper-and-page citation shown to the user.
\end{itemize}

\subsection{Data model}
Four core tables:
\begin{itemize}[leftmargin=0.7cm]
    \item \texttt{papers}(id, arxiv\_id, title, authors, year, source, created\_at)
    \item \texttt{chunks}(id, paper\_id, page, char\_start, char\_end, text, embedding \texttt{vector(384)})
    \item \texttt{queries}(id, user\_question, created\_at, latency\_ms)
    \item \texttt{answers}(id, query\_id, answer\_text, cited\_chunk\_ids, groundedness\_flag)
\end{itemize}

\subsection{CI/CD}
GitHub Actions on every push:
\begin{itemize}[leftmargin=0.7cm]
    \item Lint (ruff) and format check, then unit + integration tests (pytest).
    \item A retrieval-quality gate: run the labeled eval set and fail the build if Recall@5 drops below the current threshold.
    \item Build the Docker image and deploy to staging on a green main branch.
\end{itemize}

\section{Evaluation Criterion}
Every metric is computed automatically against a hand-labeled evaluation set of question--passage pairs, so scores are reproducible and comparable across iterations.

\subsection{Primary evaluation metric}
\textbf{Retrieval relevance --- Recall@k:} the fraction of evaluation queries for which at least one truly relevant passage appears in the top-$k$ retrieved chunks.
\begin{itemize}[leftmargin=0.7cm]
    \item Target: Recall@5 $\ge$ 0.80 on the labeled set during the pilot.
    \item Attribution: measured directly against the ground-truth labels; the same harness runs in CI, so any regression is caught on commit.
\end{itemize}

\subsection{Secondary evaluation metrics}
\begin{itemize}[leftmargin=0.7cm]
    \item \textbf{Answer groundedness / faithfulness:} percentage of answer claims that are supported by a cited chunk, checked by an LLM-as-judge pass plus a spot-checked human sample.
    \item \textbf{Mean Reciprocal Rank (MRR):} rewards ranking the correct passage higher, not just including it.
    \item \textbf{Time-to-insight reduction:} time to locate a relevant paper on the platform versus a manual keyword-search baseline.
    \item \textbf{End-to-end latency:} p95 query latency $\le$ 3 seconds, logged per query.
    \item \textbf{Reliability:} $\ge$ 99\% uptime for search and Q\&A during the pilot window.
\end{itemize}

\subsection{Pilot validation plan}
\begin{itemize}[leftmargin=0.7cm]
    \item Assemble a fixed corpus of a few hundred arXiv papers in one domain (e.g., NLP) and hand-label $\sim$50 questions with their supporting passages.
    \item Recruit 10--20 pilot users; compare time-to-insight and answer quality against the manual baseline using the labeled set and short interviews.
\end{itemize}

\section{Scalability}
\begin{enumerate}[leftmargin=0.7cm]
    \item \textbf{Scalable (theoretically):} the system absorbs growth in corpus size and concurrent queries by
    \begin{itemize}[leftmargin=0.7cm]
        \item using an HNSW index in \texttt{pgvector} for sub-linear nearest-neighbor search,
        \item decoupling ingestion (batch) from serving (online) so heavy embedding jobs never block queries,
        \item keeping the API stateless behind a load balancer.
    \end{itemize}
    \item \textbf{Operationally deployable (practically):} the system runs on common hosting via
    \begin{itemize}[leftmargin=0.7cm]
        \item a single managed Postgres instance for both metadata and vectors,
        \item a containerized app deployable to Render/Railway,
        \item a GitHub Actions pipeline for staging and production.
    \end{itemize}
\end{enumerate}

\section{Engine availability heuristic}
Every component maps to a mature, available tool, so effort goes into system design and retrieval quality rather than infrastructure:
\begin{itemize}[leftmargin=0.7cm]
    \item FastAPI for the API layer; React for the UI.
    \item \texttt{sentence-transformers} for local embeddings; Gemini 2.5 Flash for generation.
    \item PostgreSQL + \texttt{pgvector} for combined relational and vector storage.
    \item PyMuPDF for parsing; the arXiv API for ingestion.
    \item pytest for testing; GitHub Actions for CI/CD; Docker for packaging.
\end{itemize}

\section{Project scope and deliverables}
\subsection{Initial deliverable}
\begin{itemize}[leftmargin=0.7cm]
    \item Ingestion pipeline: arXiv fetch + PDF upload, parse, chunk, embed, store
    \item Semantic search endpoint returning ranked chunks with paper/page metadata
    \item Grounded Q\&A returning answers with inline source citations
    \item Labeled evaluation set and an automated harness computing Recall@k and MRR
    \item CI: lint + tests + retrieval-quality gate; CD to staging
\end{itemize}

\subsection{Subsequent deliverables}
\begin{itemize}[leftmargin=0.7cm]
    \item Per-paper and cross-paper summarization
    \item Structured knowledge extraction (task, method, dataset, metric)
    \item Hybrid retrieval (dense + BM25) with cross-encoder reranking
    \item Corpus management, filtering, and query-history views
\end{itemize}

\section{Risks and mitigations}
\begin{itemize}[leftmargin=0.7cm]
    \item \textbf{Hallucination / ungrounded answers:} strict retrieval-only prompting, mandatory citations, and an automated faithfulness check on every response.
    \item \textbf{Poor retrieval quality:} guard with the eval harness; escalate from dense to hybrid + reranking only when the data shows a gain.
    \item \textbf{PDF parsing failures:} PyMuPDF with a fallback path and per-document logging so bad files are isolated, not fatal.
    \item \textbf{Copyright / licensing:} restrict fetching to arXiv (open access) plus user-uploaded content.
    \item \textbf{Cost overrun:} embeddings run locally; the hosted LLM is called only at generation time, on a small budget.
\end{itemize}

\section{Summary}
This project delivers a deployable retrieval-augmented platform for scientific literature: arXiv and user-PDF ingestion, local embeddings with \texttt{pgvector} search, and Gemini-generated answers that cite their sources. It meets the course criteria through early time-to-value via a working vertical slice, CI/CD with an automated retrieval-quality gate for safe frequent releases, and rigorous measurable evaluation---chiefly Recall@k and answer groundedness. The engineering focus is a well-designed, cost-aware, scalable system rather than research novelty.

\end{document}